\chapter[Auto-Adapter Analysis]{Component Analysis of Auto-Adapter}
\label{chapter:autoadapter-component-analysis}

\section{Introduction}

Chapter~\ref{chapter:autoadapter-bench} evaluates how LLM backbone choice
affects driver synthesis and capability-interface use while holding the
surrounding framework fixed. This chapter changes the object of analysis from
the model to selected design choices and control boundaries in Auto-Adapter.
The system designs a reusable robot capability interface and generates a
robot-specific Python driver that implements its declared operations; the
driver is not itself a task-specific policy. This distinction separates three
kinds of evidence: producing source code, passing independent capability-level
validation, and completing a downstream task through the interface. The three
evidence blocks examine these stages without treating them as a single measure
of success.

The first study isolates the information available during driver development.
Test-driven controller synthesis can return structured \texttt{pytest}
failures, including aggregate quantities derived from simulator execution, to
guide code correction \citep{tripathi2026testdriven}. Auto-Adapter additionally
exposes a bounded, candidate-visible projection of MuJoCo state from public
development cases, such as actuator controls, joint states, and compact
trajectory summaries. The comparison holds the generation protocol and
revision opportunity fixed so that the feedback representation, rather than
additional model interaction, is the object of analysis. It asks whether
state-rich MuJoCo feedback changes synthesis outcomes beyond structured test
diagnostics; both conditions may contain simulator-derived quantities. This
public development signal remains separate from the private validation
definitions and final verdict owned by the \emph{Harness}. Related work on
simulation-based correction motivates testing simulator feedback, but does not
establish its effect in this driver-synthesis setting
\citep{kim2026physicalcorrection}.

The remaining two studies move upstream from code generation to interface and
acceptance-condition design. Task-Grounded Capability Design receives at least
twenty source-backed tasks for each robot, but no pre-authored capability
catalogue or task-to-capability mapping, and must propose five to ten capability
contracts intended for reuse across tasks. The analysis assesses whether the
resulting set remains traceable to sources, covers the source-backed task
library, and expresses semantically coherent operations; it does not treat
capability design as driver generation. The pass-criteria study examines
whether a model can select and form capability-level criteria from existing
task-level standards without weakening their source obligations. An
implementation-blind Independent Validation Compiler (IVC) audits those
criteria, binds them to private MuJoCo task instances and trusted measurement
definitions, and compiles an executable validation suite. The \emph{Harness}
retains verdict authority, so successful compilation is not itself behavioural
validation.

\section{Background}

\subsection{Physics-Grounded Feedback}

LLM coding agents use compiler and test failures as development feedback.
\citet{tripathi2026testdriven} return structured \texttt{pytest} diagnostics
that may include aggregate simulator measurements, while
\citet{kim2026physicalcorrection} use physical simulation to correct
LLM-generated tool-using primitives. The Auto-Adapter experiment asks whether
a bounded, candidate-visible projection of MuJoCo state adds information beyond
structured diagnostics. This public feedback remains distinct from the sealed
validation cases and final \emph{Harness} verdict, and simulator state does not
establish hardware behaviour.

\subsection{Capability Design}

Task-conditioned robot programming commonly generates the code or operation
needed for one instruction, potentially leaving other reusable operations
unexplored \citep{liang2023codeaspolicies,mu2024robocodex}. Task-Grounded
Capability Design instead supplies at least twenty source-backed tasks, without
a capability catalogue or task-to-capability mapping, and asks the model to
propose five to ten capabilities intended to cover the library. From the
associated task-level standards, the model also forms capability-level pass
criteria. An implementation-blind IVC audits whether those criteria preserve
source obligations, binds them to private MuJoCo task instances and trusted
measurement definitions, and compiles an executable validation suite; the
\emph{Harness} retains final verdict authority. This design is intended to give
the high-level controller a broader set of independently testable operations
to select and compose.

\section{Auto-Adapter Overview}

\section{Experiments and Results}

\subsection{MuJoCo versus \texttt{pytest} Feedback}

\subsection{Task-Grounded Capability Design}

\subsection{Pass-Criteria Design and IVC}

\subsection{Interpretation and Limitations}

\section{Summary}
